\subsection{Vue.js}
\setauthor{Zoe Emily Öllinger}
Das Frontend des Health-Monitoring-Dashboards wurde mit dem JavaScript-Framework Vue.js implementiert.
Die Entscheidung für Vue.js basiert auf seiner einfachen Integration, klaren Struktur sowie der Weiterentwicklung innerhalb der Firma.
Da innerhalb der Firma für die Entwicklung des Dashboards bereits genutzt wird, bietet uns diese Technologie perfekte Vorraussetzungen um implementiert zu werden.
Dabei bietet Vue.js uns auch eine komponentenbasierte Architektur, wodurch einzelne UI-Bestandteile unabhängig voneinander Entwickelt und wiederverwendet werden können.
Hierbei wurden uns von der Firma bereits fertige Komponenten zur Verfügung gestellt, wie zum Beispiel die Pagination.
Für uns sind die Komponenten von Vorteil, da wir im Dashboard Übersichtsseiten, Detailseiten und Statusanzeigen nutzen.
Dabei muss man beachten, dass von cer Firma für das Frontend zwei Komponenten gewünscht waren.
Einmal für die Übersicht und einmal für die Detailseite.

Ein weiterer Vorteil bei unserer Verwendung von Vue.js ist, die reaktive Datenbindung, da Änderungen an den Daten automatisch im User Interface Aktualisiert werden, wodurch eine dynamische und aktuelle Darstellung der Health-Daten gewährleitstet wird.

\cite{vuejs}


\subsubsection{Framework Background}
\setauthor{Zoe Emily Öllinger}

Vue.js wurde von Evan You entwickelt und erstmals im Jahr 2014 veröffentlicht.
Hierbei verfolgt das Framework einen progressiven Ansatz, bei dem die Entwickler selbst entscheiden können, wie umfangreich man Vue in sein Projekt integriert.

Ein zentrales Konzept von Vue ist das Reactive Data Binding.
Dabei wird das DOM automatisch aktualisiert, wenn sich Daten verändern.

Zusätzlich bietet Vue eine klare, strukturierte Trennung zwischen Template, Styling und Logik, wodurch die Wartbarkeit und Erweiterbarkeit der Anwendung besser verständlich ist.

\subsubsection{Einsatz im Frontend}
\setauthor{Zoe Emily Öllinger}

Im Rahmen dieses Projekts übernimmt Vue.js folgende zentrale Aufgaben:
\begin{compactitem}
    \item Darstellung der Health-Check-Daten
    \item Visualisierung von Statusinformationen und Fehlermeldungen
    \item Interaktion mit den APIS des Backends
    \item Dynamisches Laden und Aktualisieren von Daten
    \item Navigation zwischen verschiedenen Ansichten
\end{compactitem}
\cite{vuejs}


\subsection{UI\&UX Design mit Figma}
\setauthor{Zoe Emily Öllinger}

Für die Gestaltung der Benutzeroberfläche wurde das Design-Tool Figma verwendet.
Figma ermöglicht eine kollaborative Erstellung von UI-Designs und bietet uns umfangreiche Funktionen zur Visualisierung unserer Benutzeroberflächen.

In unserem Projekt diente Figma als zentrales Werkzeug zur Planung uns Ausarbeitung des User Interfaces und des gewünschten Workflows.
Dabei wurden von uns Layouts, Workflow und Farbkonzepte entwickelt, bevor diese dann technisch umgesetzt wurden.


\cite{figma}


\subsubsection{Designprozess}
\setauthor{Zoe Emily Öllinger}

Der Designprozess erfolgte iterativ und war stark durch regelmäßige Feedback-Stunden geprägt.
In mehreren Workflow und Feedback-Sessions wurden die Entwürfe gemeinsam besprochen und anschließend verbessert und an die Anforderungen der Mitarbeiter angepasst.

Diese Vorgehensweise ermöglichte es uns, frühzeitig Verbesserungen zu erkennen und entsprechend darauf zu reagieren und die Benutzerfreundlichkeit auf Wunsch der Mitarbeiter kontinuierlich zu optimieren.

\cite{figma}


\subsubsection{Vorteile von Figma}
\setauthor{Zoe Emily Öllinger}

Der Einsatz von Figma bringt mehrere Vorteile mit sich:
\begin{compactitem}
    \item Zentrale und kollaborative Designplattform
    \item Schnelle Erstellung von Prototypen
    \item Einfache Abstimmung im Team
    \item Möglichkeit zur direkten Visualisierung von User Flows
    \item Konsistente Gestaltung durch wiederverwendbare Komponenten
\end{compactitem}


\cite{figma}

\subsection{Benutzerfreundlichkeit und Interaktion}
\setauthor{Zoe Emily Öllinger}

Ein zentrales Ziel des Frontends war die Entwicklung einer intuitiven und übersichtlichen Benutzeroberflache.

Besonderes Augenmerk wurde auf folgende Aspekte gelegt:
\begin{compactitem}
    \item Übersichtliche Darstellung von Statusinformationen
    \item Klare visuelle Hervorhebung von Fehlern und Warnungen
    \item Schnelle Navigation zwischen verschiedenen Bereichen
    \item Minimierung von Ladezeiten durch effiziente Datenabfragen
    \item Kommentarsektion zur Plattform-Übergreifenden bearbeitung
    \item Anzeige von gravierenden Fehlern mit Verlinkung zur Fehlerbehebung und Dokumentation
\end{compactitem}


\newpage